% Bloque independiente. Sustituir los campos por evidencia de su proyecto.
\bloque{\personaUno}
\section{RPC}
\begin{frame}{RPC: una operación en otro proceso}\relax
{\normalsize Un cliente solicita la ejecución de un procedimiento en un servidor y recibe su resultado.\par}
\vspace{1mm}
\begin{columns}[T,onlytextwidth]
\begin{column}{.46\textwidth}
\textbf{Contrato compartido}\par
{\small Operación, argumentos y respuesta. El \emph{stub} representa al servicio.\par}
\vspace{1mm}
\textbf{Serialización}\par
{\small Convierte los datos en un mensaje; el receptor reconstruye sus valores.\par}
\end{column}
\begin{column}{.50\textwidth}
\centering
\begin{tikzpicture}[font=\small]
\node[nodo,minimum width=4.8cm,minimum height=.75cm] (c) {Cliente y \emph{stub}};
\node[nodo,minimum width=4.8cm,minimum height=.75cm,below=1.15cm of c] (s) {Servidor: ejecuta la operación};
\draw[-{Stealth},thick,draw=Azul] ([xshift=-1.25cm]c.south) -- node[left,font=\scriptsize]{Petición} ([xshift=-1.25cm]s.north);
\draw[-{Stealth},thick,draw=Azul] ([xshift=1.25cm]s.north) -- node[right,font=\scriptsize]{Respuesta} ([xshift=1.25cm]c.south);
\end{tikzpicture}
\end{column}
\end{columns}
\vspace{1mm}
{\small\textbf{La red importa:} hay latencia y posibles fallos. El formato y el transporte dependen de la implementación.\par}
\vspace{1.5mm}
{\scriptsize Contrato y stub \cite{zhang2023rpc,arcuri2025evomaster}; procesamiento y latencia \cite{seemakhupt2023}.}
\note{Explicar RPC como modelo de invocación remota. Recorrer contrato, stub, serialización, petición, ejecución y respuesta. Aclarar que RPC no exige HTTP/2 ni Protocol Buffers. La guía personal de teoría está en personal/belinda/README.md y se excluye localmente de Git.}
\end{frame}

\begin{frame}{Ejemplo sencillo}\relax
\textbf{Operación}\quad \campo{Operación pequeña con JSON-RPC u otra implementación RPC}\par
\medskip
\textbf{Entrada}\quad \campo{Dato de prueba}\par
\textbf{Salida}\quad \campo{Resultado realmente obtenido}\par
{\footnotesize Contrato y entradas de una API RPC \cite{arcuri2025evomaster}.}
\guia{Mostrar método, argumentos, respuesta y los procesos que intervienen.}
\end{frame}

\begin{frame}{SGRA: consulta académica como servicio remoto}\relax
{\small\textbf{Propuesta futura:} convertir la consulta Java interna en una llamada a otro proceso.\par}
\vspace{0mm}
\centering
% Diagrama propio de Belinda Toaquiza, derivado del bosquejo del 19-11-2026.
% Vista de la consulta seleccionada; propuesta futura, no despliegue verificado.
\begin{tikzpicture}[
  x=1cm,y=1cm,
  bloque/.style={draw=Azul!60,fill=Bruma,rounded corners=2pt,
    align=center,text width=2.6cm,minimum height=1.08cm,inner sep=5pt,font=\fontsize{9}{10.5}\selectfont},
  flujo/.style={-{Stealth[length=2mm]},thick,draw=Azul},
  vuelta/.style={-{Stealth[length=2mm]},thick,draw=Gris},
  rotulo/.style={font=\fontsize{7.5}{9}\selectfont,align=center,text=Tinta}
]
\node[bloque,text width=1.8cm] (angular) at (0,0) {\textbf{Angular}\\Estudiante};
\node[bloque] (backend) at (3.7,0) {\textbf{Backend SGRA}\\Refuerzos\\Cliente JSON-RPC};
\node[bloque] (academico) at (8.8,0) {\textbf{Servicio académico}\\Valida identidad\\Consulta JDBC};
\draw[draw=Azul!65,dashed,rounded corners=3pt]
  (2.12,-1.15) rectangle (5.28,1.05);
\draw[draw=Azul!65,dashed,rounded corners=3pt]
  (7.22,-1.15) rectangle (10.38,1.05);
\node[rotulo,anchor=south] at (3.7,1.10) {Proceso del backend};
\node[rotulo,anchor=south] at (8.8,1.10) {Proceso separado propuesto};
\draw[flujo] ([yshift=2mm]angular.east) -- node[above,rotulo]{HTTP(S)} ([yshift=2mm]backend.west);
\draw[vuelta] ([yshift=-2mm]backend.west) -- ([yshift=-2mm]angular.east);
\node[rotulo,anchor=north] at (1.7,-.5) {Respuesta\\para la pantalla};
\draw[flujo] ([yshift=2.5mm]backend.east) -- node[above,rotulo]{Usuario + asignatura} ([yshift=2.5mm]academico.west);
\draw[vuelta] ([yshift=-2.5mm]academico.west) -- node[below,rotulo]{Docente o\\sin asignación} ([yshift=-2.5mm]backend.east);
\node[rotulo,text=Azul,anchor=south] at (6.25,1.12) {\textbf{JSON-RPC 2.0}\\sobre HTTP(S)};
\node[bloque,text width=2.6cm,minimum height=.7cm,fill=white] (bd) at (8.8,-2.25) {\textbf{PostgreSQL}\\Base existente};
\draw[flujo] ([xshift=-3.5mm]academico.south) -- node[left,rotulo]{SQL / JDBC} ([xshift=-3.5mm]bd.north);
\draw[vuelta] ([xshift=3.5mm]bd.north) -- node[right,rotulo]{Resultado} ([xshift=3.5mm]academico.south);
\node[rotulo,anchor=north west,text width=6.8cm,align=left] at (-1.0,-1.5) {
  \textbf{Operación propuesta}\\
  \texttt{consultarDocenteAsignado}\\
  Entrada: usuario autenticado y asignatura.\\
  El contexto de identidad debe validarse entre servicios.
};
\end{tikzpicture}

\par\vspace{0mm}
\raggedright
{\scriptsize\textbf{Supuesto:} base existente sin réplicas. \textbf{Base del análisis:} \href{https://github.com/KhrizGomez/back-SGRA/blob/e28bcfac58b92d44904939dbfd95ae9a75777c5b/src/main/java/com/CLMTZ/Backend/service/reinforcement/student/impl/StudentRequestServiceImpl.java\#L41}{código de SGRA}, revisado el 19-11-2026.}
\note{Diagrama propio derivado del bosquejo de Belinda del 19-11-2026. En el flujo actual, el controlador y StudentRequestServiceImpl usan StudentCatalogService dentro del backend. La propuesta reemplaza esa dependencia por una llamada JSON-RPC a otro proceso. La identidad del usuario proviene del contexto autenticado; debe validarse entre servicios. PostgreSQL responde al servicio, no directamente al navegador. La separación conserva inicialmente la base existente como supuesto; no se infiere replicación ni una etiqueta CAP. La foto original queda en evidencias/Belinda T/esquema-mano/sgra-bosquejo-preliminar-consulta-docente-belinda-2026-11-19.jpeg.}
\end{frame}

\begin{frame}{Decisiones al distribuir el proyecto}\relax
\small
\begin{tabularx}{\linewidth}{@{}lX@{}}
\toprule
Comunicación & \campo{Operación, contrato y llamada bloqueante o no bloqueante. Justificar.}\\
Transparencias & \campo{Una ofrecida y una limitada, con ejemplo para el usuario.}\\
CAP & \campo{Topología supuesta, enlace que se corta y conducta de lectura/escritura.}\\
PACELC & \campo{Con la red sana, qué consistencia y latencia buscaríamos y cómo medirlas.}\\
\bottomrule
\end{tabularx}
\guia{Vincular cada decisión con el flujo del diagrama. Precisar réplicas y garantías antes de asignar una etiqueta CAP.}
\end{frame}

\begin{frame}{Fallo hipotético y lección para el proyecto}\relax
\small
\textbf{Escenario supuesto}\quad \campo{Durante una operación, el servicio o enlace falla.}\par
\textbf{Causa y recuperación}\quad \campo{Cadena causal, detección y respuesta propuesta.}\par
{\footnotesize Diagnosticar la latencia de llamadas \cite{traini2023delag}.}\par
\textbf{Comparación}\quad \campo{Sistema real comparable, fuente y compromiso técnico.}\par
\textbf{Lección}\quad \campo{Qué implementaríamos primero y con qué prueba lo validaríamos.}\par
\guia{Separar el resultado esperado de una prueba ya realizada. Si prueban el fallo, registrar pasos, salida y fecha como simulación.}
\end{frame}
